\chapter{The Weil conjectures: exact statement and conventions}
\label{chap:statement}

\section{Arithmetic data attached to a variety}

Let \(k_0=\Fq\), let \(k=\Fbar\), and let \(X_0\) be a scheme of finite
type over \(k_0\).  For every \(n\geq1\), the set
\(X_0(\mathbf F_{q^n})\) is finite.  We write
\[
  N_n(X_0)=\#X_0(\mathbf F_{q^n}).
\]
The Hasse--Weil zeta function packages all these point counts into one
formal series.

\begin{definition}
\label{def:hasse-weil-zeta}
The zeta function of \(X_0\) is
\[
 Z(X_0,T)
 =\exp\left(\sum_{n\geq1}N_n(X_0)\frac{T^n}{n}\right)
 \in1+T\mathbf Q[[T]].
\]
\end{definition}

The denominators \(n\) make it nonobvious from this definition that the
coefficients are integers.  Integrality becomes transparent after
rewriting the series as an Euler product.  If \(x\) is a closed point,
its residue field is finite over \(k_0\), and we set
\[
  \deg(x)=[\kappa(x):k_0].
\]
The closed-point formula, proved without cohomology in
\cref{prop:euler-product}, is
\[
  Z(X_0,T)=\prod_{x\in|X_0|}
       (1-T^{\deg(x)})^{-1}\in1+T\mathbf Z[[T]].
  \tag{1.1}\label{eq:zeta-euler-intro}
\]
This is a formal product: only finitely many factors can affect a fixed
coefficient.  Analytic convergence is neither assumed nor needed.

\section{Frobenius and Tate twists}

Put \(X=X_0\times_{k_0}k\).  The symbol
\[
 F\colon X\longrightarrow X
\]
always denotes geometric Frobenius, inverse to the \(q\)-power arithmetic
Frobenius.  On geometric points it generates the same finite orbits as
arithmetic Frobenius, but in the opposite direction.  On etale
cohomology, taking the inverse changes every eigenvalue, so the
distinction is indispensable.

\begin{convention}
Our Tate-twist convention is
\[
 F\mid\Ql(1)=q^{-1},\qquad
 F\mid\Ql(-1)=q.
\]
Thus \(F\) acts by \(q^d\) on \(\Ql(-d)\).  If \(X_0\) is smooth and
proper of dimension \(d\), the trace map identifies
\[
 H^{2d}_{\et}(X,\Ql)\simeq\Ql(-d)
\]
when \(X\) is connected.  The top reciprocal Frobenius root is therefore
\(q^d\), not \(q^{-d}\).
\end{convention}

Fix a prime \(\ell\neq\operatorname{char}(k_0)\).  For a smooth proper
\(X_0\), define
\[
 V_{i,\ell}(X_0)=H^i_{\et}(X,\Ql),\qquad
 P_{i,\ell}(T)=\det(1-TF\mid V_{i,\ell}(X_0)).
 \tag{1.2}\label{eq:def-Pi}
\]
These spaces are finite-dimensional and vanish outside
\(0\leq i\leq2d\).  If, over \(\Qbarl\),
\[
 P_{i,\ell}(T)=\prod_{j=1}^{b_{i,\ell}}
       (1-\alpha_{ij}T),
\]
then the \(\alpha_{ij}\) are called the reciprocal roots, or Frobenius
eigenvalues.  The ordinary roots of the polynomial are
\(\alpha_{ij}^{-1}\).

\section{Weights: the initial and final formulations}

There are two logically different formulations of purity.  The
sheaf-theoretic theorem first says that an eigenvalue
\(\alpha\in\Qbarl^\times\) has
\[
 |\iota(\alpha)|=q^{w/2}
 \tag{1.3}\label{eq:universal-iota-weight}
\]
for every abstract field isomorphism
\(\iota\colon\Qbarl\xrightarrow{\sim}\mathbf C\).  It does not initially
assert that \(\alpha\) is algebraic over \(\mathbf Q\).  After
no-cancellation and rational factor recovery are proved, \(\alpha\) is
known to be an algebraic number, and \eqref{eq:universal-iota-weight}
becomes the classical assertion
\[
 |\tau(\alpha)|=q^{w/2}
\]
for every embedding
\(\tau\colon\mathbf Q(\alpha)\hookrightarrow\mathbf C\).

\begin{warning}
It is not enough to choose one isomorphism
\(\Qbarl\simeq\mathbf C\).  The final rationality of each weight factor
uses stability under \(\operatorname{Aut}(\Qbarl/\mathbf Q)\).  That
stability follows because \eqref{eq:universal-iota-weight} is quantified
over every \(\iota\): composing \(\iota\) with an automorphism preserves
the same absolute value.  This quantifier is a mathematical input, not a
matter of exposition.
\end{warning}

\section{The classical theorem}

\begin{theorem}[The Weil conjectures]
\label{thm:weil-conjectures-statement}
Let \(X_0\) be a smooth, projective, geometrically connected variety of
pure dimension \(d\) over \(\Fq\).  Then:

\begin{enumerate}
\item \emph{Cohomological rationality.}  For every auxiliary prime
      \(\ell\neq\operatorname{char}(\Fq)\),
      \[
        Z(X_0,T)
        =\prod_{i=0}^{2d}P_{i,\ell}(T)^{(-1)^{i+1}}.
        \tag{1.4}\label{eq:weil-cohomological-factorization}
      \]
      In particular \(Z(X_0,T)\) is a rational function.

\item \emph{Integral, realization-independent factors.}  There are
      polynomials \(P_i(T)\in\mathbf Z[T]\), independent of \(\ell\),
      such that \(P_i=P_{i,\ell}\) for every auxiliary prime.  They have
      constant term \(1\), and
      \[
        P_0(T)=1-T,\qquad P_{2d}(T)=1-q^dT.
        \tag{1.5}\label{eq:weil-endpoints}
      \]
      Consequently
      \[
        Z(X_0,T)=
        \frac{P_1(T)P_3(T)\cdots}
        {(1-T)P_2(T)P_4(T)\cdots(1-q^dT)},
      \]
      with only degrees in the interval \([0,2d]\) included.

\item \emph{Riemann hypothesis.}  Every reciprocal root \(\alpha\) of
      \(P_i\) is an algebraic integer, and for every embedding
      \(\tau\colon\mathbf Q(\alpha)\hookrightarrow\mathbf C\),
      \[
        |\tau(\alpha)|=q^{i/2}.
        \tag{1.6}\label{eq:weil-rh}
      \]
      Equivalently, every ordinary root of \(P_i(T)\) has absolute value
      \(q^{-i/2}\) under every complex embedding.

\item \emph{Betti numbers.}  The integer
      \[
        b_i=\deg P_i=\dim_{\Ql}H^i_{\et}(X,\Ql)
        \tag{1.7}\label{eq:weil-betti}
      \]
      is independent of \(\ell\).  Poincare duality gives
      \(b_i=b_{2d-i}\), and geometric connectedness gives
      \(b_0=b_{2d}=1\).  If \(X_0\) is obtained by good reduction from a
      smooth proper variety in characteristic zero, then \(b_i\) equals
      the \(i\)-th singular Betti number of the associated compact
      complex manifold.

\item \emph{Functional equation.}  If
      \[
        \chi=\sum_{i=0}^{2d}(-1)^ib_i,
      \]
      then a sign \(\varepsilon\in\{1,-1\}\) exists such that
      \[
        Z(X_0,q^{-d}T^{-1})
        =\varepsilon q^{d\chi/2}T^\chi Z(X_0,T).
        \tag{1.8}\label{eq:weil-functional-equation}
      \]
      If \(D_i=\det(F\mid H^i_{\et}(X,\Ql))\), then
      \[
        \varepsilon=(-1)^\chi q^{-d\chi/2}
          \prod_{i=0}^{2d}D_i^{(-1)^i}.
        \tag{1.9}\label{eq:epsilon-determinant}
      \]
      Poincare duality proves that the right side is a sign.
\end{enumerate}
\end{theorem}

\section{What the Betti-number assertion means}

In characteristic zero, a smooth projective complex variety has
topological cohomology groups whose dimensions are its classical Betti
numbers.  An arbitrary variety over a finite field need not possess a
lift to characteristic zero, so one cannot define its \(b_i\) by first
choosing such a lift.  The unconditional definition in
\eqref{eq:weil-betti} is etale.  The substantive assertion is that the
dimension is the same for every \(\ell\).  When a smooth proper lift
exists, smooth proper base change and the etale--singular comparison
theorem prove agreement with topology.

The theorem does not assert hard Lefschetz as part of its statement.
Hard Lefschetz gives further structure on the \(b_i\), but it is not
needed in the Laumon route used in these notes.  Poincare duality alone
is enough for the symmetry \(b_i=b_{2d-i}\) and for the functional
equation.

\section{The geometric-connectedness caveat}

For any smooth proper \(X_0\),
\[
 H^0_{\et}(X,\Ql)
\]
is the vector space of functions on the finite set \(\pi_0(X)\), and
\(F\) permutes this set.  Hence in general
\[
 P_{0,\ell}(T)=\det(1-TF\mid\Ql^{\pi_0(X)}).
 \tag{1.10}\label{eq:general-P0}
\]
It is \(1-T\) precisely when \(X\) is connected.  Duality makes the top
factor the same permutation determinant with each eigenvalue multiplied
by \(q^d\).

For instance, \(\Spec\mathbf F_{q^m}\) is connected over \(\Fq\), but
its geometric fiber consists of \(m\) points cyclically permuted by
Frobenius.  Its \(P_0\) is \(1-T^m\).  Thus ``connected'' cannot replace
``geometrically connected'' in the endpoint assertion.  The
cohomological trace formula, purity, and factor arguments remain true
without this hypothesis; only the simplified endpoints
\eqref{eq:weil-endpoints} must be replaced by
\eqref{eq:general-P0} and its dual.

For a smooth scheme, geometric connectedness implies geometric
irreducibility.  Indeed, a regular scheme has disjoint irreducible
components: if two components met, the local ring at an intersection
point would have more than one minimal prime and would not be a domain,
whereas a regular local ring is a domain.  The components are therefore
open and closed.  This explains why older statements often say
``nonsingular projective absolutely irreducible variety.''

\section{Logical order of the proof}

The proof does not establish the five clauses of
\cref{thm:weil-conjectures-statement} independently.
It follows this directed chain:
\[
\begin{array}{c}
\text{trace formula}+\text{finite-dimensional cohomology}\\
\Downarrow\\
\text{\(\Ql\)-rational determinant expression for }Z\\
\Downarrow\\
\text{Hankel descent gives }Z\in\mathbf Q(T),\\[1mm]
\text{while Poincare duality gives the functional equation};
\end{array}
\]
\[
\begin{array}{c}
\text{Fourier--Deligne transform}+\text{curve estimate}\\
\Downarrow\\
\text{Deligne--Laumon direct-image theorem}\\
\Downarrow\\
\text{smooth-proper universal realization-purity}\\
\Downarrow\\
\text{no cancellation between weights}\\
\Downarrow\\
\text{algebraicity, rational factors, integral factors,
and \(\ell\)-independence}.
\end{array}
\]
In particular, the individual \(P_{i,\ell}\) are not assumed integral,
rational over \(\mathbf Q\), or \(\ell\)-independent during the proof of
purity.  Their alternating product is known earlier, but it is purity
which makes the factors recoverable from that product.

The route is exclusively the Laumon--Deligne Weil II route.  We do not
use the Weil I chain through Lefschetz pencils, Picard--Lefschetz
monodromy, Chebotarev, or tensor-power amplification.  Keeping one route
prevents deep inputs from being imported twice under different names and
makes the dependency audit suitable for formalization.

\begin{formalization}
A reusable theorem statement should separate the following structures:
\begin{itemize}
\item a counting sequence \(n\mapsto N_n\) and a locally finite
      collection of closed-point degrees;
\item a bounded graded family of finite-dimensional Frobenius modules;
\item a trace-formula equality connecting the first two;
\item universal realization-purity, arithmetic algebraicity, and
      algebraic integrality as three different predicates;
\item a perfect graded similitude pairing with multiplier \(q^d\);
\item geometric connectedness as the hypothesis which simplifies the
      endpoint permutation representation.
\end{itemize}
The polynomial \(P_i=\det(1-TF_i)\) should be constructed from the
Frobenius module rather than stored as redundant data.  The final
\(\ell\)-independent polynomial should be returned by a uniqueness
theorem for weight factors of the integral zeta function.
\end{formalization}
